\documentclass[11pt]{article}
\usepackage[includeheadfoot, top=1.0in, bottom=1.0in, hmargin=1.0in]{geometry}
\usepackage[utf8]{inputenc}
\usepackage{fancyhdr}
\usepackage{url}
\pagestyle{fancy}
\usepackage{setspace}
\usepackage{tabularx}
\usepackage{xcolor}
\usepackage{cancel}
\usepackage{amsmath,amsfonts}

\lhead{Astronomy Lab II}
\rhead{Spring 2023}
\lfoot{Mead \& Golant}
\rfoot{Tue 7-10pm}
\cfoot{\thepage}

\begin{document}

\begin{center}
\huge{Lab 1: Orders of Magnitude}\\ \medskip \Large{January 24, 2023}
\end{center}

%%%%%%%%%%%%%%%%%%%%%%% INTRODUCTION %%%%%%%%%%%%%%%%%%%%%%%

\section{Introduction: An Astronomical Undertaking}

The Earth has a mass of about 6 \emph{septillion} kilograms -- that's 6 followed by \emph{24} zeros. The Sun's mass is about 300,000 times that. And a typical galaxy weighs in at about 100 \emph{billion} Suns. Better yet, a large galaxy cluster can contain up to \emph{1,000} galaxies. That's a lot of mass.

\medskip \noindent
It's not a coincidence that, in colloquial English, the word ``astronomical'' is synonymous with ``\emph{very very large}.'' In astronomy, we study an extremely wide range of scales, from the tiniest electron to the largest cluster of galaxies. As such, we need some type of mathematical framework for dealing with very big (and very small) numbers -- this framework is called \textbf{scientific notation} (a remarkably generic name). 

\medskip \noindent
In this lab, you will learn how to carry out astronomical calculations using \emph{scientific notation}, how to properly represent astronomical measurements with \emph{units}, and how to understand the vastness of the Universe in terms of \emph{orders of magnitude}. This lab will provide you with the basic quantitative skills required to be an amateur astronomer.

\bigskip


%%%%%%%%%%%%%%%%%%%%%%% PRACTICE %%%%%%%%%%%%%%%%%%%%%%%
\section{Let's Practice!}

Solve the following problems using what we've just discussed about \textit{orders of magnitude, scientific notation, significant figures,} and \textit{units}.  \textbf{Record your solutions in your lab notebook.}  Be sure to show your work for full credit!
\begin{enumerate}
    \item \textcolor{blue}{6 pts} Light travels at a speed of $1.8 \times 10^7$ kilometers per minute. The average distance from the Earth to the Sun is $1.5 \times 10^8$ kilometers. Find the time (in minutes) that it takes light to travel from the Sun to the Earth. Give your answer in scientific notation with the proper number of significant figures. If the Sun were to suddenly stop emitting light, how long would it take us (on Earth) to notice this? \textcolor{red}{$8.3$ minutes}
    
    \item \textcolor{blue}{6 pts} Cosmic voids -- the emptiest regions in the Universe -- have densities lower than $8 \times 10^{-28}$ kilograms per cubic meter. Voids can also be extremely large. Taking the volume of a void to be $4 \times 10^{56}$ cubic meters, find the total amount of mass (in kilograms) contained in a typical void. Express your answer in scientific notation with the proper number of significant figures. On Earth, the density of dry air at sea level is 1.2 kilograms per cubic meter. By how many orders of magnitude is the density of air greater than the density in cosmic voids? \textcolor{red}{$3.2 \times 10^{29}$ kg; 27 or 28 OoM}
    
    \item \textcolor{blue}{6 pts} The speed at which the Earth orbits around the Sun depends on the \emph{sum} of the Earth's mass and the Sun's mass. The mass of the Earth is $5.972 \times 10^{24} \, {\rm kg}$ and the mass of the Sun is $1.989 \times 10^{30} \, {\rm kg}$. Add these two masses together and express your answer in scientific notation with 2 significant figures. Does the mass of the Earth play a significant role in determining the Earth's orbital speed around the Sun? \textcolor{red}{$1.989 \times 10^{30}$ kg; NO}

    \item \textcolor{blue}{9 pts} The Hubble constant, which tells us how fast the Universe is expanding, is commonly written with the weird units ``kilometers per second per megaparsec'' (or $\frac{\rm km/s}{\rm Mpc}$; 1 Mpc = $10^6$ pc). Convert $70 \, \frac{\rm km/s}{\rm Mpc}$ to units of $\frac{\rm km/s}{\rm km}$ (what units is this equivalent to?). You can obtain a rough estimate for the age of the Universe (in seconds) by taking the \emph{reciprocal} of your answer (i.e., by dividing 1 by your answer). What do you get for the age of the Universe? What is this value in \emph{years}, rather than in seconds? \textcolor{red}{$2.269 \times 10^{-18} s^{-1}$; $4.407 \times 10^{17} s$; $13.97 \times 10^9 yrs$}
\end{enumerate}


%%%%%%%%%%%%%%%%%%%%%%% SCALING THE UNIVERSE %%%%%%%%%%%%%%%%%%%%%%%
\section{Scaling the Universe}

Let's apply what we've learned so far to construct a \emph{scale model} of the Universe. Let's scale the diameter of the Sun ($1.39 \times 10^6$ kilometers) down to the size of a ballpoint pen tip (1.00 millimeter) to get a sense of just how large the Universe really is. \textbf{Record your solutions in your lab write-up}.


\begin{enumerate}
    \item \textcolor{blue}{3 pts} First, let's set up the \emph{scale factor}, $F$, that'll connect our scale model to the real thing. The tip of a ballpoint pen is $F$ times \emph{smaller} than the Sun, so $$R_{\rm pen} = F \times R_{\rm sun}.$$ As such, find $F$ (to three significant figures) by dividing the size of the pen tip by the size of the Sun. Make sure that your units are consistent when dividing! (1 km = $10^6$ mm) \textcolor{red}{$F = 7.19 \times 10^{-13}$}
    \item Now, scale down the following distances by multiplying by the scale factor, $F$.
    \begin{enumerate}
        \item \textcolor{blue}{3 pts} The distance from the Sun to the Earth (1.0 AU) \textcolor{red}{$7.19 \times 10^{-13} \, \rm{AU}$}
        \item \textcolor{blue}{3 pts} The distance from the Sun to Jupiter (5.2 AU) \textcolor{red}{$3.74 \times 10^{-12} \, \rm{AU}$}
        \item \textcolor{blue}{3 pts} The distance to the edge of the Solar System (100,000 AU) \textcolor{red}{$7.19 \times 10^{-8} \, \rm{AU}$}
        \item \textcolor{blue}{3 pts} The distance to the nearest star (4.25 ly) \textcolor{red}{$3.06 \times 10^{-12} \, \rm{ly}$}
        \item \textcolor{blue}{3 pts} The distance to the center of the Milky Way (8 kiloparsecs) \textcolor{red}{$5.76 \times 10^{-12} \, \rm{kpc}$}
        \item \textcolor{blue}{3 pts} The distance to the Andromeda Galaxy (0.89 megaparsecs) \textcolor{red}{$6.40 \times 10^{-13} \, \rm{Mpc}$}
        \item \textcolor{blue}{3 pts} The distance to the edge of the observable Universe ($4.65\times 10^{10}$ ly) \textcolor{red}{$3.34 \times 10^{-2} \, \rm{ly}$}
    \end{enumerate}
\end{enumerate}

\bigskip

\section{Fermi Estimation}
Throughout this lab, we've talked a good bit about significant figures and the \emph{precision} of measurements. But, in many problems, all we really need is a rough estimate of the answer -- the order of magnitude alone can tell us a fair amount about the problem at hand (as we've discussed, comparing orders of magnitude is a quick and easy way of getting a feel for how big or how small a quantity is). In the process of \textbf{\emph{Fermi estimation}} (also, aptly, called \emph{order-of-magnitude} estimation), we seek to estimate (usually to only one significant figure) the answer to a problem by using rough, intuitive arithmetic rather than fully rigorous calculations. For instance, if we wanted to Fermi estimate the number of jelly beans in a jar, we shouldn't need to rigorously compute the volume of each bean and the spacing between each bean -- we can make some simplifying assumptions about the shape of the beans (i.e., we can estimate the beans as being spherical) and about the geometry of the bean pile (i.e., we can estimate the beans as being in a simple lattice pattern) and use these assumptions to obtain a rough order-of-magnitude estimate for the number of beans.

\medskip \noindent
Fermi estimation is a difficult skill to perfect, but it's a very important tool for developing quantitative intuition. So, let's try a few examples. \textbf{Record your solutions in your lab write-up}. \textit{Show your work and be explicit about your assumptions.} Don't worry about your estimations being exactly right -- the thought process is what counts.
\begin{enumerate}
    \item \textcolor{blue}{6 pts} In an act of rebellion against your poor lab instructor, you declare scientific notation to be useless and decide to write out $10^{\rm one \, billion}$ (that's 1 followed by one \emph{billion} zeros) in its entirety. Assuming you're using a standard word processor with 12 pt font and 1 inch margins, estimate how many sheets of $8.5'' \times 11''$ paper you'd need to fully write out $10^{\rm one \, billion}$. Give your answer to one significant figure. If you placed all these sheets of paper in a single stack, approximately how tall (in meters, to one significant figure) would that stack be? 
    
    \item \textcolor{blue}{6 pts} The main Columbia campus extends from 114th street to 120th street, between Broadway and Amsterdam Avenue. Assuming there are no buildings or other obstructions in this rectangle, approximately how many people could fit into Columbia's main campus (without needing to stack people on top of one other)? 
    
    \item \textcolor{blue}{10 pts} Read a little bit about the \textbf{Drake equation} online and explain how this could be framed as a Fermi estimation problem. Using the Drake equation, provide an order-of-magnitude answer to this Fermi problem (you may use numbers from Wikipedia or from other websites). 
    
    \item \textcolor{blue}{10 pts} Come up with your own Fermi estimation problem and write out a short solution!
\end{enumerate}

%%%%%%%%%%%%%%%%%%%%%%% CONCLUSIONS %%%%%%%%%%%%%%%%%%%%%%%
\section{Wrapping things up}

Complete this section by yourself. \textbf{Record your responses in your lab write-up}. 
\begin{enumerate}
    \item \textcolor{blue}{4 pts} Why is scientific notation useful in astronomy (and in quantitative disciplines in general)?
    
    \item \textcolor{blue}{4 pts} Describe qualitatively what an ``order of magnitude'' is. Why do we care about orders of magnitude?
    
    \item \textcolor{blue}{4 pts} Why is it important to include units when you report a measurement? When might one want to perform a unit conversion?
    
    \item What is one remaining question you have after this lab?  I will try to answer remaining question at the beginning of next lab. 
    
    \item Please provide some feedback on how today's lab session was run/organized. What worked well? What could be done to make future lab sessions better? What do you hope to get out of this class, and how can the instructors make this happen?
\end{enumerate}

\end{document}

